\documentclass[12pt]{article}
\usepackage[utf8]{inputenc}
\usepackage[T1]{fontenc}
\usepackage{hyperref}
\title{Kumbaran Language Resources for K-12 Curriculum Development}
\author{Compiled by Claude Sonnet 5.5 (claude-sonnet-5-5) for Andr\'as Kornai}
\date{}
\begin{document}
\maketitle

\section{Introduction}
Kumbaran (ISO 639-3 wkb; Glottocode kumb1267) is a Dravidian language of India with about 10,000 speakers as of 2004, spoken by a Scheduled Caste community in the Wayanad area of Kerala. Its classification is unsettled: Wikipedia treats it as unclassified Southern Dravidian and notes it may be Malayalamoid. The survey figure of 48\% cognacy with Telugu comes from an unpublished 2004 report and is weak evidence. The family label ``Bookkeeping'' in the task header appears to be a Glottolog bookkeeping category, and the reference description gives Dravidian.

\section{Primary Educational Resources}
None found. The language is not known to be taught in schools. Malayalam is the likely school language in the region.

\section{Linguistic and Reference Materials}
\begin{itemize}
\item Overview: \url{https://en.wikipedia.org/wiki/Kumbaran_language}
\item Glottolog reference entry for Jose Mathew (2004), the survey cited for the language: \url{https://glottolog.org/resource/reference/id/151540}
\item Wikipedia's article on Malayalamoid languages discusses possible relatedness: \url{https://en.wikipedia.org/wiki/Malayalamoid_languages}
\end{itemize}
Other results (Kurumba, Kurux, Kumbarti) concern different languages and were not used.

\section{Religious and Community Resources}
No scripture or community organisation for Kumbaran was found.

\section{Digital and Audio Resources}
No audio, video, corpus or software resources were found.

\section{Curriculum Development Context}
Almost nothing is documented. There is no grammar, dictionary, orthography or teaching material in the search results. A curriculum would first require descriptive fieldwork and community involvement.

\end{document}